\documentclass[11pt,a4paper]{article}

% --- Packages ---
\usepackage[margin=.75in]{geometry}
\usepackage[utf8]{inputenc}
\usepackage[T1]{fontenc}
\usepackage{lmodern} % Fix for font not found error
\usepackage{titlesec}
\usepackage{enumitem}
\usepackage{hyperref}
\usepackage{xcolor}
\usepackage{fancyhdr}
\usepackage{comment}
\usepackage{cleveref}
\usepackage[bottom]{footmisc}
\usepackage{graphicx}
\usepackage{longtable}
\usepackage{booktabs}
\usepackage{array}
\usepackage{calc}
\usepackage{etoolbox}

% Standard Pandoc setup
\providecommand{\tightlist}{%
  \setlength{\itemsep}{0pt}\setlength{\parskip}{0pt}}

\crefformat{section}{\S#2#1#3}
\crefformat{subsection}{\S#2#1#3}
\crefformat{subsubsection}{\S#2#1#3}
\crefformat{figure}{#2Figure~#1#3}
\crefformat{footnote}{#2\footnotemark[#1]#3}

\linespread{0.9} % Riduce lo spazio tra le linee

% --- Colors ---
\definecolor{primary}{RGB}{0, 0, 0}
\definecolor{links}{HTML}{0000EE}
\definecolor{blue}{HTML}{26428B}

\hypersetup{
    colorlinks=true,
    linkcolor=links,
    urlcolor=links,
    citecolor=links,
}

% --- Section Formatting ---
\titleformat{\section}
  {\normalfont\Large\bfseries\color{blue}}{\thesection}{1em}{}
\titleformat{\subsection}
{\normalfont\large\bfseries\color{blue}}{\thesubsection}{1em}{}
\titleformat{\subsubsection}
{\normalfont\bfseries\color{blue}}{\thesubsubsection}{1em}{}

% --- Header and Footer ---
\pagestyle{fancy}
\fancyhead[L]{\small Giuseppe Capasso}
\fancyhead[R]{\small vapt}
\fancyfoot[C]{\thepage}
\renewcommand{\headrulewidth}{0.4pt}

% --- Custom Title Command ---
\newcommand{\proposaltitle}[1]{
    \begin{center}
        \vspace*{0.1cm}
        {\LARGE \bfseries \color{blue} #1} \\[1.5ex]
        {\large \textbf{Giuseppe Capasso}} \\[1ex]
                \small{
            \vspace{0.1cm}
            \textbf{Materia:} vapt\\
        }
            \end{center}
    \vspace{0.3cm}
    \hrule height 0.5pt
    \vspace{0.5cm}
}

\begin{document}

\proposaltitle{Introduzione}

\section{Introduzione}\label{introduzione}

In questa lezione, ci focalizzeremo sugli attacchi destinati
all'infrastruttura di rete. In particolare, copriremo attacchi alle LAN
e ai principali protocolli di livello 3 dello stack TCP/IP.

\section{Fondamenti di Rete}\label{fondamenti-di-rete}

Per comprendere come un\textquotesingle infrastruttura possa essere
compromessa, è indispensabile padroneggiare il suo funzionamento in
condizioni di normalità. In questo capitolo ripasseremo i concetti
architetturali alla base delle comunicazioni locali e geografiche,
focalizzandoci sui protocolli e sui meccanismi che governano il Data
Link Layer e il Network Layer del modello OSI.

\subsubsection{Identità in Rete: Indirizzamento Fisico e
Logico}\label{identituxe0-in-rete-indirizzamento-fisico-e-logico}

La comunicazione di rete si basa su un sistema di doppio indirizzamento,
necessario per gestire sia la consegna locale
(all\textquotesingle interno dello stesso segmento) sia
l\textquotesingle instradamento globale (attraverso reti diverse).

\begin{itemize}
\item
  \textbf{MAC Address (Media Access Control):
  L\textquotesingle Indirizzo Fisico (Layer 2)}

  \begin{itemize}
  \item
    È un identificatore univoco a 48 bit, tipicamente espresso in
    formato esadecimale (es. 00:1A:2B:3C:4D:5E).
  \item
    È "bruciato" (burned-in) nella scheda di rete (NIC) dal produttore e
    opera esclusivamente a livello di rete locale (LAN).
  \item
    I primi 24 bit rappresentano l\textquotesingle OUI (Organizationally
    Unique Identifier), che identifica il vendor, mentre i restanti 24
    bit sono assegnati specificamente a quella singola interfaccia.
  \end{itemize}
\item
  \textbf{IP Address (Internet Protocol): L\textquotesingle Indirizzo
  Logico (Layer 3)}

  \begin{itemize}
  \item
    Nel caso dell\textquotesingle IPv4, è un indirizzo gerarchico a 32
    bit, diviso in quattro ottetti decimali (es. 192.168.1.10).
  \item
    A differenza del MAC, l\textquotesingle IP non è legato
    all\textquotesingle hardware ma è assegnato logicamente
    (staticamente o dinamicamente via DHCP) e definisce sia la rete di
    appartenenza (Network ID) sia lo specifico host (Host ID).
  \item
    È essenziale per il routing: permette ai router di inoltrare i
    pacchetti tra reti geograficamente o logicamente separate.
  \end{itemize}
\end{itemize}

\subsubsection{La Meccanica dello
Switch}\label{la-meccanica-dello-switch}

Lo switch è il cuore nevralgico della LAN moderna. A differenza del
vecchio hub, che replicava ogni segnale elettrico su tutte le porte
(causando collisioni e inefficienze), lo switch è un apparato
intelligente che inoltra i frame Layer 2 in modo selettivo.

Il suo funzionamento si basa sulla \textbf{MAC Address Table} (o CAM
Table):

\begin{enumerate}
\def\labelenumi{\arabic{enumi}.}
\item
  \textbf{Apprendimento (Learning):} Quando uno switch riceve un frame,
  legge il MAC address \emph{sorgente} e lo associa alla porta fisica da
  cui è entrato, salvando questa informazione nella sua tabella.
\item
  \textbf{Inoltro (Forwarding):} Lo switch consulta poi il MAC address
  \emph{destinazione} del frame. Se questo indirizzo è presente nella
  sua tabella, inoltra il frame esclusivamente alla porta associata.
\item
  \textbf{Flooding:} Se il MAC destinazione non è presente nella tabella
  (Unknown Unicast), oppure se si tratta di un indirizzo di Broadcast
  (es. FF:FF:FF:FF:FF:FF), lo switch inoltra il frame su tutte le porte
  attive, ad eccezione di quella da cui lo ha ricevuto.
\end{enumerate}

\subsubsection{Il Protocollo ARP: Il Ponte tra Logico e
Fisico}\label{il-protocollo-arp-il-ponte-tra-logico-e-fisico}

Affinché un host A possa inviare un pacchetto IP a un host B sulla
stessa rete locale, deve prima incapsulare il pacchetto in un frame
Ethernet. Per farlo, ha bisogno di conoscere il MAC address di B. Qui
entra in gioco l\textquotesingle{}\textbf{ARP (Address Resolution
Protocol)}.

Il processo di risoluzione avviene in due fasi:

\begin{enumerate}
\def\labelenumi{\arabic{enumi}.}
\item
  \textbf{ARP Request:} L\textquotesingle host A invia un messaggio in
  broadcast (FF:FF:FF:FF:FF:FF) a tutta la sottorete chiedendo:
  \emph{"Chi ha l\textquotesingle indirizzo IP 192.168.1.50? Risponda
  indicando il proprio MAC address"}.
\item
  \textbf{ARP Reply:} L\textquotesingle host B (e solo lui) riconosce il
  proprio IP, genera una risposta unicast diretta
  all\textquotesingle host A comunicando il proprio MAC address.
\end{enumerate}

Una volta ricevuta la risposta, l\textquotesingle host A salva
l\textquotesingle associazione IP-MAC nella propria \textbf{ARP Cache}
locale, in modo da non dover ripetere la richiesta per le comunicazioni
successive a breve termine.

\subsubsection{La Vulnerabilità di
ARP}\label{la-vulnerabilituxe0-di-arp}

Se dal punto di vista operativo l\textquotesingle ARP è un capolavoro di
efficienza, dal punto di vista della sicurezza è un protocollo
intrinsecamente vulnerabile. È stato progettato con
un\textquotesingle architettura \textbf{stateless} (senza memoria di
stato) e si basa sul principio della totale "fiducia implicita"
(implicit trust).

Le due debolezze architetturali fondamentali che un Penetration Tester
sfrutta sono:

\begin{enumerate}
\def\labelenumi{\arabic{enumi}.}
\item
  \textbf{Mancanza di Autenticazione:} Non esiste alcun meccanismo
  crittografico per verificare che l\textquotesingle host che invia un
  messaggio ARP sia effettivamente il legittimo proprietario di
  quell\textquotesingle indirizzo IP.
\item
  \textbf{Accettazione Acritica:} Un host aggiornerà la propria cache
  ARP anche se riceve una \emph{ARP Reply} non sollecitata (ovvero,
  senza aver mai inviato una precedente \emph{ARP Request}).
\end{enumerate}

Questa dinamica viene spesso innescata abusando dei \textbf{Gratuitous
ARP}, ovvero messaggi di broadcast legittimi usati tipicamente dagli
apparati di rete per annunciare un cambio di indirizzo MAC (es. in
scenari di failover come l\textquotesingle HSRP) o per rilevare
conflitti IP.

\subsubsection{ARP Spoofing (o ARP Poisoning):
L\textquotesingle Attacco}\label{arp-spoofing-o-arp-poisoning-lattacco}

L\textquotesingle attacco di \textbf{ARP Spoofing} consiste
nell\textquotesingle inviare pacchetti ARP forgiati (falsificati) sulla
rete locale per "avvelenare" (poisoning) la cache ARP degli host target.

\textbf{La dinamica dell\textquotesingle attacco:}
L\textquotesingle attaccante vuole intercettare il traffico tra una
Vittima (es. 192.168.1.50) e il Default Gateway (es. 192.168.1.1).

\begin{enumerate}
\def\labelenumi{\arabic{enumi}.}
\item
  L\textquotesingle attaccante invia una \emph{ARP Reply} falsificata
  alla Vittima dichiarando: \emph{"Il MAC address
  dell\textquotesingle IP 192.168.1.1 (Gateway) è il mio MAC address"}.
\item
  Parallelamente, l\textquotesingle attaccante invia una \emph{ARP
  Reply} falsificata al Gateway dichiarando: \emph{"Il MAC address
  dell\textquotesingle IP 192.168.1.50 (Vittima) è il mio MAC address"}.
\end{enumerate}

\textbf{Gli impatti principali:}

\begin{itemize}
\item
  \textbf{Man-in-the-Middle (MitM):} Tutto il traffico in uscita dalla
  vittima e diretto verso internet passerà prima attraverso la macchina
  dell\textquotesingle attaccante (che agirà da router invisibile
  facendo IP forwarding). Questo permette lo \emph{sniffing} di
  credenziali in chiaro, il dirottamento di sessioni e la manipolazione
  dei pacchetti in transito.
\item
  \textbf{Denial of Service (DoS):} Se l\textquotesingle attaccante
  associa l\textquotesingle IP del gateway a un MAC address inesistente
  (o semplicemente non inoltra i pacchetti intercettati), la vittima
  subisce un \emph{blackholing} del traffico, perdendo totalmente la
  connettività esterna.
\end{itemize}

\subsubsection{Contromisure e Mitigazioni
Architetturali}\label{contromisure-e-mitigazioni-architetturali}

In un ambiente enterprise o durante un\textquotesingle attività di
remediation post-VA/PT, la sicurezza del Layer 2 non può essere affidata
al caso. Le difese si dividono in operative e infrastrutturali:

\begin{itemize}
\item
  \textbf{Tabelle ARP Statiche (Mitigazione Operativa):} Consiste nel
  forzare manualmente l\textquotesingle associazione IP-MAC nei sistemi
  operativi degli host critici o sul router (arp -s su Windows/Linux). È
  una soluzione invulnerabile allo spoofing, ma assolutamente
  \textbf{non scalabile} e inadatta a reti dinamiche con client DHCP.
  Viene usata solo in casi limite (es. tra due server ad altissima
  criticità).
\item
  \textbf{Dynamic ARP Inspection - DAI (Mitigazione Infrastrutturale):}
  È lo standard de facto per la difesa a livello enterprise. Il DAI è
  una feature di sicurezza configurabile sugli switch (come i Cisco
  Catalyst o la serie 8000v) che intercetta tutti i pacchetti ARP in
  transito sulle porte non fidate (untrusted).

  \begin{itemize}
  \item
    \textbf{Come funziona:} Lo switch analizza ogni pacchetto ARP e lo
    scarta se rileva che il MAC address e l\textquotesingle IP address
    non corrispondono a un\textquotesingle associazione legittima.
  \item
    \textbf{Da dove prende le informazioni valide?} Il DAI lavora in
    tandem con il \textbf{DHCP Snooping Binding Database}, una tabella
    costruita dallo switch osservando le normali transazioni DHCP. Se un
    host ha regolarmente ottenuto un IP dal server DHCP, lo switch sa
    esattamente quale MAC address è autorizzato a usare
    quell\textquotesingle IP. Se arriva un pacchetto ARP anomalo, viene
    bloccato, l\textquotesingle evento viene loggato e la porta può
    essere messa in stato di \emph{err-disable}.
  \end{itemize}
\end{itemize}

\subsubsection{Il Sistema DNS: La Risoluzione dei
Nomi}\label{il-sistema-dns-la-risoluzione-dei-nomi}

Mentre macchine e router ragionano per indirizzi IP e MAC, gli esseri
umani utilizzano nomi mnemonici (es. www.azienda.local). Il
\textbf{Domain Name System (DNS)} è l\textquotesingle infrastruttura
globale e gerarchica (Layer 7) che si occupa di questa traduzione.

\begin{itemize}
\item
  \textbf{Funzionamento:} Quando un client deve raggiungere un dominio,
  il suo resolver locale invia una query (tipicamente UDP sulla porta
  53) al server DNS configurato (spesso fornito dal DHCP).
\item
  \textbf{Risoluzione:} Se il server DNS locale possiede
  l\textquotesingle informazione nella sua cache o nella sua zona
  autoritativa, risponde immediatamente con l\textquotesingle indirizzo
  IP corrispondente. Altrimenti, interroga ricorsivamente altri server
  DNS (Root servers, TLD servers, Authoritative servers) fino a ottenere
  la risposta da inoltrare al client.
\item
  Il corretto funzionamento del DNS è vitale: senza di esso, sebbene la
  connettività di base (Layer 3) sia intatta, la navigazione e
  l\textquotesingle accesso ai servizi risultano di fatto impossibili
  per un utente standard.
\end{itemize}

\subsubsection{Ripasso Teorico: Il Protocollo
DHCP}\label{ripasso-teorico-il-protocollo-dhcp}

Il \textbf{Dynamic Host Configuration Protocol (DHCP)} è un protocollo
di tipo client-server che automatizza l\textquotesingle assegnazione
degli indirizzi IP e di altri parametri di rete critici. Senza DHCP,
ogni dispositivo dovrebbe essere configurato manualmente (Static IP),
rendendo la gestione di reti moderne virtualmente impossibile.

\paragraph{Funzionamento e Porte}\label{funzionamento-e-porte}

Il DHCP opera a livello applicativo (Layer 7) ma interagisce
direttamente con i livelli inferiori per la consegna dei pacchetti
quando l\textquotesingle host non ha ancora un\textquotesingle identità
logica.

\begin{itemize}
\item
  \textbf{Protocollo di Trasporto:} UDP.
\item
  \textbf{Porta Server:} 67 (ascolta le richieste dei client).
\item
  \textbf{Porta Client:} 68 (riceve le risposte dal server).
\end{itemize}

\paragraph{Il Processo DORA}\label{il-processo-dora}

La transazione standard tra un client e un server DHCP avviene in
quattro fasi principali, riassunte dall\textquotesingle acronimo
\textbf{DORA}:

\begin{enumerate}
\def\labelenumi{\arabic{enumi}.}
\item
  \textbf{D}iscover (Client -\textgreater{} Server):
  L\textquotesingle host appena connesso invia un pacchetto in broadcast
  (255.255.255.255) per individuare i server DHCP disponibili nella rete
  locale.
\item
  \textbf{O}ffer (Server -\textgreater{} Client): Uno o più server
  rispondono proponendo una configurazione (IP, Subnet, Gateway). A
  questo livello, l\textquotesingle IP non è ancora assegnato, è solo
  "prenotato".
\item
  \textbf{R}equest (Client -\textgreater{} Server): Il client accetta
  formalmente una delle offerte ricevute (solitamente la prima) inviando
  un altro broadcast. Questo serve anche a informare gli altri eventuali
  server DHCP che la loro offerta non è stata scelta.
\item
  \textbf{A}cknowledge (Server -\textgreater{} Client): Il server
  conferma l\textquotesingle assegnazione e invia i dettagli finali. Da
  questo momento l\textquotesingle host è ufficialmente configurato.
\end{enumerate}

\paragraph{Il Concetto di Lease
(Locazione)}\label{il-concetto-di-lease-locazione}

L\textquotesingle indirizzo IP non viene regalato al client, ma
"affittato" per un tempo determinato, chiamato \textbf{Lease Time}.

\begin{itemize}
\item
  \textbf{T1 (Renewal):} Al 50\% della durata del lease, il client tenta
  di rinnovarlo contattando il server in unicast.
\item
  \textbf{T2 (Rebinding):} Se il server non risponde,
  all\textquotesingle87.5\% della durata il client invia un broadcast
  per cercare \emph{qualsiasi} server DHCP che possa estendere il lease.
\item
  \textbf{Expiration:} Se il tempo scade senza rinnovo,
  l\textquotesingle host perde l\textquotesingle IP e deve ricominciare
  il processo dal Discover.
\end{itemize}

\paragraph{DHCP Relay Agent}\label{dhcp-relay-agent}

In reti enterprise segmentate, il server DHCP spesso non risiede nella
stessa VLAN dei client. Poiché i broadcast non superano i router, si
utilizza il \textbf{DHCP Relay} (comando Cisco ip helper-address). Il
router riceve il broadcast del client, lo trasforma in un pacchetto
unicast e lo instrada verso l\textquotesingle IP del server DHCP remoto.

\begin{itemize}
\tightlist
\item
  \emph{Nota per il VA/PT:} Il Relay Agent aggiunge informazioni al
  pacchetto (come l\textquotesingle indirizzo
  dell\textquotesingle interfaccia sorgente), permettendo al server di
  capire da quale pool deve pescare l\textquotesingle IP.
\end{itemize}

\subsubsection{First Hop Redundancy}\label{first-hop-redundancy}

In un\textquotesingle infrastruttura di rete tradizionale, il Default
Gateway rappresenta un critico \emph{single point of failure} (punto
singolo di guasto). Se l\textquotesingle interfaccia del router che
funge da gateway si disconnette o subisce un disservizio,
l\textquotesingle intero segmento di rete locale perde la capacità di
comunicare con l\textquotesingle esterno, indipendentemente dalla
stabilità del resto dell\textquotesingle infrastruttura.

Per mitigare questo rischio ingegneristico sono stati sviluppati i
protocolli della famiglia FHRP (First Hop Redundancy Protocol). Tra
questi, l\textquotesingle{}\textbf{HSRP (Hot Standby Router Protocol)} è
lo standard proprietario Cisco più diffuso. L\textquotesingle obiettivo
dell\textquotesingle HSRP è garantire un failover trasparente e
immediato, raggruppando più router fisici all\textquotesingle interno di
un unico cluster logico.

\subsubsection{Architettura e Funzionamento
dell\textquotesingle HSRP}\label{architettura-e-funzionamento-dellhsrp}

In una topologia HSRP, due o più router collaborano per presentarsi agli
host della LAN come un singolo "Virtual Router". Questo apparato
virtuale è definito da due elementi fondamentali:

\begin{itemize}
\item
  \textbf{Virtual IP (VIP):} Un indirizzo IP logico condiviso tra i
  router del gruppo. Questo è l\textquotesingle indirizzo che viene
  configurato come default gateway sui client (staticamente o tramite
  DHCP).
\item
  \textbf{Virtual MAC Address:} Un indirizzo fisico generato
  algoritmicamente e associato al VIP. Nelle implementazioni classiche
  (HSRPv1), ha il formato 0000.0c07.acXX, dove gli ultimi due caratteri
  esadecimali rappresentano l\textquotesingle ID del gruppo HSRP.
\end{itemize}

\textbf{Il Processo Elettivo:} Affinché non ci siano conflitti di
instradamento, i router del gruppo eleggono un \textbf{Active Router},
che si assume l\textquotesingle onere esclusivo di inoltrare il traffico
destinato al Virtual MAC, e uno \textbf{Standby Router}, che rimane in
ascolto pronto a subentrare in caso di guasto.

L\textquotesingle elezione si basa su un valore di \textbf{Priority} (da
0 a 255, con 100 come valore di default). Il router con la priorità più
alta vince. Per mantenere l\textquotesingle accordo sullo stato del
cluster, i router si scambiano periodicamente pacchetti di \emph{Hello}
in multicast (all\textquotesingle indirizzo 224.0.0.2 su porta UDP
1985). Se lo Standby Router smette di ricevere i pacchetti Hello
dall\textquotesingle Active Router per un tempo superiore
all\textquotesingle{}\emph{Hold Time} (di default 10 secondi), dichiara
il leader caduto e assume immediatamente il ruolo attivo.

\subsubsection{Vulnerabilità HSRP: Il Takeover del
Gateway}\label{vulnerabilituxe0-hsrp-il-takeover-del-gateway}

Esattamente come per l\textquotesingle ARP, l\textquotesingle HSRP
nativo è stato concepito in un\textquotesingle era di totale fiducia
infrastrutturale. Nelle sue configurazioni di base,
l\textquotesingle HSRP è privo di autenticazione o utilizza, al massimo,
una debole autenticazione in chiaro (spesso lasciata sul default
"cisco").

Questa debolezza espone la rete a un attacco critico di \textbf{HSRP
Takeover}. Un Penetration Tester (o un attaccante) posizionato nello
stesso dominio di broadcast può sfruttare le dinamiche del protocollo
per diventare lui stesso il Default Gateway della rete.

\textbf{La dinamica dell\textquotesingle attacco:}

\begin{enumerate}
\def\labelenumi{\arabic{enumi}.}
\item
  \textbf{Reconnaissance:} L\textquotesingle attaccante si mette in
  ascolto passivo sulla rete catturando i pacchetti multicast 224.0.0.2.
  Ispezionando gli \emph{Hello} packet dell\textquotesingle HSRP,
  estrapola il Virtual IP, l\textquotesingle ID del gruppo,
  l\textquotesingle indirizzo IP del router attivo, la sua priorità e
  l\textquotesingle eventuale password in chiaro.
\item
  \textbf{Injection (Takeover):} Utilizzando tool di network injection
  (come \emph{Yersinia} o script Scapy personalizzati),
  l\textquotesingle attaccante inizia a trasmettere messaggi di
  \emph{Hello} contraffatti sulla rete, dichiarando di appartenere allo
  stesso gruppo HSRP, conoscendo la password corretta, ma impostando una
  \textbf{Priority massima (255)}.
\item
  \textbf{Preemption:} Se sui router legittimi è abilitata la funzione
  di \emph{Preemption} (una feature che permette a un router con
  priorità superiore di spodestare immediatamente il leader attuale
  senza aspettare che cada), la macchina dell\textquotesingle attaccante
  viene eletta istantaneamente nuovo Active Router.
\end{enumerate}

\textbf{Gli impatti principali:} A questo punto, tutti i client della
rete invieranno il loro traffico in uscita alla macchina
dell\textquotesingle attaccante. Questo scenario permette un attacco
Man-in-the-Middle perfetto su scala di rete, consentendo
l\textquotesingle intercettazione, l\textquotesingle alterazione o il
blocco (DoS) del traffico verso internet o verso altre VLAN.

\subsubsection{Contromisure: Blindare la
Ridondanza}\label{contromisure-blindare-la-ridondanza}

La mitigazione degli attacchi contro i protocolli di routing e
ridondanza richiede di eliminare la "fiducia cieca" e imporre una
validazione rigorosa dei messaggi.

\begin{itemize}
\item
  \textbf{Autenticazione Crittografica:} La soluzione definitiva è
  abbandonare l\textquotesingle autenticazione in chiaro (o la sua
  assenza) e configurare l\textquotesingle autenticazione basata su
  \textbf{MD5} (o algoritmi superiori se supportati) per il gruppo HSRP.
  In questo modo, ogni pacchetto \emph{Hello} include un hash
  crittografico basato su una chiave segreta condivisa tra i soli router
  legittimi. Eventuali pacchetti iniettati da un attaccante verranno
  immediatamente scartati dal cluster, poiché
  l\textquotesingle attaccante non possiede la chiave per generare
  l\textquotesingle hash corretto.
\item
  \textbf{Filtri Multicast e ACL:} Aggiungere Access Control List sulle
  interfacce di switch e router per limitare quali porte fisiche possono
  originare traffico multicast destinato all\textquotesingle indirizzo
  224.0.0.2
\end{itemize}

Ottima scelta strategica. Lavorare a coppie sui Cisco 2960 fisici con le
VM in bridge abbatte il carico computazionale e simula un environment
"bare metal" molto più realistico rispetto alle reti virtuali chiuse.

Per mantenere la continuità con la stesura precedente, nominerò questa
sezione "Capitolo 3", dato che il Capitolo 2 lo abbiamo dedicato alla
teoria e agli attacchi HSRP.

Ecco la dispensa per il laboratorio.

\section{Laboratorio Pratico - ARP Spoofing, Phishing e Mitigazioni
Layer
2}\label{laboratorio-pratico---arp-spoofing-phishing-e-mitigazioni-layer-2}

Questo laboratorio trasforma i concetti teorici di vulnerabilità del
Layer 2 in un attacco reale. Lavorerete in coppie: uno studente assumerà
il ruolo dell\textquotesingle attaccante (Red Team),
l\textquotesingle altro quello della vittima. Successivamente,
consoliderete l\textquotesingle infrastruttura implementando le
contromisure sugli switch Cisco 2960.

\subsubsection{Setup dell\textquotesingle Ambiente: VirtualBox in
"Bridged Mode"}\label{setup-dellambiente-virtualbox-in-bridged-mode}

Di default, VirtualBox configura le interfacce di rete in modalità NAT,
creando una rete isolata dove la VM esce su internet "dietro"
l\textquotesingle host fisico. Per gli attacchi Layer 2 (come
l\textquotesingle ARP Spoofing), l\textquotesingle attaccante deve
risiedere nello \textbf{stesso dominio di broadcast} della vittima. È
quindi necessario "ponticellare" la scheda di rete virtuale di Kali con
quella fisica del vostro PC, connessa allo switch Cisco.

\textbf{Procedura su VirtualBox:}

\begin{enumerate}
\def\labelenumi{\arabic{enumi}.}
\item
  Spegnete la macchina virtuale Kali Linux.
\item
  Aprite le \textbf{Impostazioni} della VM e andate nella sezione
  \textbf{Rete}.
\item
  Alla voce \textbf{Connesso a}, selezionate dal menu a tendina
  \textbf{Scheda con bridge} (Bridged Adapter).
\item
  Alla voce \textbf{Nome}, selezionate la scheda di rete Ethernet fisica
  del vostro host (quella fisicamente cablata allo switch 2960).
\item
  Espandete la sezione \textbf{Avanzate} e, alla voce \textbf{Modalità
  promiscua}, selezionate \textbf{Permetti a tutte} (Allow All). Questo
  è cruciale per permettere alla scheda di catturare il traffico non
  destinato direttamente al suo MAC address durante lo sniffing.
\item
  Avviate Kali e verificate di aver ottenuto un IP dalla rete del
  laboratorio (ip a).
\end{enumerate}

\subsubsection{Fase di Attacco: Man-in-the-Middle e Credential
Harvesting}\label{fase-di-attacco-man-in-the-middle-e-credential-harvesting}

In questa fase, l\textquotesingle attaccante intercetterà il traffico
della vittima e clonerà una pagina di login per sottrarre le
credenziali.

\paragraph{Step 1: Abilitare l\textquotesingle IP
Forwarding}\label{step-1-abilitare-lip-forwarding}

Prima di "avvelenare" la rete, l\textquotesingle attaccante deve
istruire il proprio kernel Linux a instradare i pacchetti. Se non lo si
fa, il traffico della vittima verrebbe droppato dalla macchina
dell\textquotesingle attaccante, causando un Denial of Service (la
vittima non navigherebbe più) e rivelando l\textquotesingle attacco.

Da terminale su Kali (come root):

Bash

echo 1 \textgreater{} /proc/sys/net/ipv4/ip\_forward

\paragraph{Step 2: Esecuzione dell\textquotesingle ARP
Spoofing}\label{step-2-esecuzione-dellarp-spoofing}

Utilizzeremo la suite dsniff. Assumiamo che:

\begin{itemize}
\item
  eth0 = interfaccia di Kali
\item
  192.168.1.50 = IP della Vittima
\item
  192.168.1.254 = IP del Default Gateway
\end{itemize}

L\textquotesingle attaccante deve aprire due finestre di terminale per
ingannare entrambe le direzioni:

\textbf{Terminale 1 (Inganna la Vittima):}

Bash

arpspoof -i eth0 -t 192.168.1.50 192.168.1.254

\textbf{Terminale 2 (Inganna il Gateway):}

Bash

arpspoof -i eth0 -t 192.168.1.254 192.168.1.50

Ora tutto il traffico tra la vittima e il gateway passa fisicamente
attraverso l\textquotesingle interfaccia
dell\textquotesingle attaccante.

\paragraph{Step 3: Clonazione del portale con SE Toolkit
(SET)}\label{step-3-clonazione-del-portale-con-se-toolkit-set}

Sfrutteremo il Social-Engineer Toolkit per clonare la pagina di login di
Twitter e catturare le credenziali immesse dalla vittima.

\begin{enumerate}
\def\labelenumi{\arabic{enumi}.}
\item
  Avviate il tool da terminale: setoolkit
\item
  Navigate il menu interattivo:

  \begin{itemize}
  \item
    Selezionate 1 (Social-Engineering Attacks)
  \item
    Selezionate 2 (Website Attack Vectors)
  \item
    Selezionate 3 (Credential Harvester Attack Method)
  \item
    Selezionate 2 (Site Cloner)
  \end{itemize}
\item
  \textbf{IP address for the POST back:} Inserite l\textquotesingle IP
  della vostra macchina Kali.
\item
  \textbf{Enter the url to clone:} Inserite http://twitter.com (o la
  pagina di login esatta).
\end{enumerate}

\emph{Nota Tecnica: Nel web moderno, domini come Twitter implementano
HSTS (HTTP Strict Transport Security), che forza la crittografia HTTPS e
rende i browser restii ad accettare certificati non validi o connessioni
declassate. Ai fini del laboratorio, se la vittima naviga direttamente
sull\textquotesingle IP dell\textquotesingle attaccante (o tramite una
regola di DNS Spoofing associata per un dominio http-only), vedrà la
pagina clonata e i log di SET registreranno la password in chiaro non
appena la vittima tenterà il login.}

L\textquotesingle ARP Spoofing ci permette di ricevere i pacchetti della
vittima, ma il \textbf{DNS Spoofing} ci permette di decidere dove la
vittima debba andare quando digita un URL (es. www.google.it). In questo
scenario, noi agiremo come un server DNS malevolo che fornisce risposte
falsificate.

\paragraph{1. Preparazione dell\textquotesingle Attaccante
(Kali)}\label{preparazione-dellattaccante-kali}

Per prima cosa, dobbiamo assicurarci che le richieste DNS della vittima
non raggiungano mai il server DNS legittimo (come quello di Google
8.8.8.8). Usiamo \textbf{IPtables} per intercettare e bloccare il
traffico DNS (porta 53 UDP) che la vittima tenta di inoltrare:

Bash

\# Blocca il traffico DNS in uscita dalla vittima che stiamo inoltrando

sudo iptables -A FORWARD -p udp -\/-dport 53 -j DROP

Successivamente, creiamo un file di configurazione (spesso chiamato
hosts.txt) che indichi al nostro tool di spoofing (come dnschef o
ettercap) quale dominio dirottare verso il nostro IP:

\begin{quote}
\textbf{Esempio file hosts.txt:} \textless ip-attaccante\textgreater{}
www.target-domain.com 10.10.10.11
\href{http://www.google.it/}{www.google.it}
\end{quote}

\paragraph{\texorpdfstring{\textbf{2. Esecuzione
dell\textquotesingle Attacco}}{2. Esecuzione dell\textquotesingle Attacco}}\label{esecuzione-dellattacco}

Lanciamo dnsspoof indicando l\textquotesingle interfaccia di rete e il
file dei nomi creato:

Bash

sudo dnsspoof -i eth0 -f fake\_hosts.txt

\paragraph{3. Verifica Lato Vittima
(Windows/Linux)}\label{verifica-lato-vittima-windowslinux}

Per testare se l\textquotesingle attacco sta funzionando, dobbiamo agire
sulla macchina vittima. È fondamentale svuotare la cache DNS locale,
altrimenti il computer userà l\textquotesingle indirizzo IP corretto
salvato in precedenza senza chiedere al server (noi).

Su \textbf{Windows}, apri il prompt dei comandi e digita:

\begin{enumerate}
\def\labelenumi{\arabic{enumi}.}
\item
  \textbf{Svuota la cache:} ipconfig /flushdns
\item
  \textbf{Interroga il dominio:} nslookup www.google.it
\end{enumerate}

\textbf{Risultato atteso:} Se l\textquotesingle attacco ha successo, il
comando nslookup restituirà l\textquotesingle indirizzo IP della
macchina Kali dell\textquotesingle attaccante invece
dell\textquotesingle IP reale di Google. Da questo momento, ogni
tentativo di navigazione verso quel dominio porterà la vittima su un
server web controllato dall\textquotesingle attaccante (es. una pagina
di phishing).

\subsubsection{Approfondimento: HTTPS e HSTS (Le difese
moderne)}\label{approfondimento-https-e-hsts-le-difese-moderne}

Anche se il DNS Spoofing ha successo, l\textquotesingle attaccante si
scontrerà quasi sempre con due tecnologie che proteggono
l\textquotesingle integrità del web.

\paragraph{\texorpdfstring{\textbf{HTTPS (HyperText Transfer Protocol
Secure)}}{HTTPS (HyperText Transfer Protocol Secure)}}\label{https-hypertext-transfer-protocol-secure}

HTTPS cifra il traffico e utilizza \textbf{certificati SSL/TLS} per
garantire l\textquotesingle identità del sito.

\begin{itemize}
\item
  \textbf{Cosa succede:} Se dirottiamo la vittima che cerca
  www.google.it sul nostro server web (Kali), il suo browser chiederà un
  certificato valido per quel dominio.
\item
  \textbf{L\textquotesingle ostacolo:} Dato che non possediamo il
  certificato firmato da una CA (Certificate Authority) per Google, il
  browser mostrerà un enorme avviso di sicurezza (\textbf{"La
  connessione non è privata"}). Se la vittima non forza manualmente
  l\textquotesingle accesso, l\textquotesingle attacco fallisce.
\end{itemize}

\paragraph{\texorpdfstring{\textbf{HSTS (HTTP Strict Transport
Security)}}{HSTS (HTTP Strict Transport Security)}}\label{hsts-http-strict-transport-security}

L\textquotesingle HSTS è un meccanismo di policy di sicurezza che
istruisce il browser a comunicare con un sito \textbf{esclusivamente}
tramite HTTPS.

\begin{itemize}
\item
  \textbf{Il funzionamento:} La prima volta che un utente visita un sito
  (es. Facebook), lo switch riceve un header HSTS. Da quel momento, il
  browser ricorderà per mesi (o anni) che quel sito non deve mai essere
  caricato in HTTP semplice.
\item
  \textbf{Impatto sul DNS Spoofing:} Anche se
  l\textquotesingle attaccante prova a fare "SSL Stripping" (forzare il
  downgrade a HTTP), il browser si rifiuterà di caricare la pagina non
  cifrata. Inoltre, in presenza di HSTS, gli errori di certificato sono
  \textbf{non ignorabili}: l\textquotesingle utente non avrà nemmeno il
  tasto "Procedi comunque".
\end{itemize}

\subsubsection{Fase di Difesa: Hardening del Cisco
2960}\label{fase-di-difesa-hardening-del-cisco-2960}

Ora che avete compreso quanto sia banale compromettere una rete
"piatta", passiamo alla mitigazione. Vi collegherete in console al Cisco
2960 per implementare due livelli di sicurezza fondamentali.

\paragraph{Difesa 1: Port Security}\label{difesa-1-port-security}

Il Port Security previene il MAC Flooding e limita i movimenti laterali
di un attaccante, vincolando una porta a uno specifico MAC address.

Accedete alla configurazione dell\textquotesingle interfaccia a cui è
collegata la vittima (es. FastEthernet 0/1):

Cisco CLI

Switch\textgreater{} enable

Switch\# configure terminal

Switch(config)\# interface fa0/1

Switch(config-if)\# switchport mode access

Switch(config-if)\# switchport port-security

Switch(config-if)\# switchport port-security maximum 1

Switch(config-if)\# switchport port-security mac-address sticky

Switch(config-if)\# switchport port-security violation restrict

\begin{itemize}
\item
  \textbf{sticky:} Lo switch apprende dinamicamente il primo MAC address
  che fa traffico sulla porta e lo salva in running-config come
  legittimo.
\item
  \textbf{restrict:} Se l\textquotesingle attaccante prova a fare
  spoofing usando un MAC address diverso (o se si collega un altro PC
  alla stessa porta), lo switch droppa i pacchetti illeciti e genera un
  log SNMP/Syslog senza spegnere la porta (a differenza
  dell\textquotesingle opzione shutdown).
\end{itemize}

\paragraph{Difesa 2: Segmentazione tramite
VLAN}\label{difesa-2-segmentazione-tramite-vlan}

L\textquotesingle ARP Spoofing funziona solo se attaccante e vittima
condividono lo stesso dominio di broadcast. Metterli in VLAN separate
distrugge alla base l\textquotesingle efficacia di questo attacco Layer
2.

Creiamo due VLAN e assegniamo le porte:

Cisco CLI

Switch(config)\# vlan 10

Switch(config-vlan)\# name IT\_ADMIN

Switch(config-vlan)\# exit

Switch(config)\# vlan 20

Switch(config-vlan)\# name GUEST

Switch(config-vlan)\# exit

Switch(config)\# interface fa0/1

Switch(config-if)\# switchport access vlan 10

Switch(config-if)\# exit

Switch(config)\# interface fa0/2

Switch(config-if)\# switchport access vlan 20

Switch(config-if)\# end

\paragraph{Step di Verifica}\label{step-di-verifica}

Per assicurarvi che le difese siano attive e funzionanti, utilizzate i
comandi di verifica (show) dal privilegio exec:

\begin{enumerate}
\def\labelenumi{\arabic{enumi}.}
\item
  \textbf{Verifica Port Security:}

  Cisco CLI

  Switch\# show port-security interface fa0/1

  \emph{Analizzate l\textquotesingle output: Verificate che lo stato sia
  Secure-up, che il Maximum MAC Addresses sia 1 e osservate il counter
  dei Security Violation dopo aver provato a forzare un cambio di MAC
  address da Kali.}
\item
  \textbf{Verifica VLAN:}

  Cisco CLI

  Switch\# show vlan brief

  \emph{Assicuratevi che le porte fa0/1 e fa0/2 siano isolate nelle
  rispettive VLAN.}
\item
  \textbf{Il Test Definitivo:} Provate a lanciare nuovamente
  l\textquotesingle attacco arpspoof da Kali. Dato che i broadcast ARP
  (FF:FF:FF:FF:FF:FF) non superano i confini della VLAN, le \emph{ARP
  Request/Reply} dell\textquotesingle attaccante non raggiungeranno mai
  la vittima, neutralizzando del tutto l\textquotesingle attacco.
\end{enumerate}

\subsubsection{Rilevamento e Difesa Lato Host (Sulla Macchina
Vittima)}\label{rilevamento-e-difesa-lato-host-sulla-macchina-vittima}

Se l\textquotesingle infrastruttura di rete non è protetta da misure
come la Dynamic ARP Inspection, la macchina vittima può comunque
utilizzare i comandi del proprio sistema operativo per rilevare un
attacco di Man-in-the-Middle in corso e per immunizzarsi.

\paragraph{Step 1: L\textquotesingle Indicatore di Compromissione
(Rilevamento)}\label{step-1-lindicatore-di-compromissione-rilevamento}

Prima che il Red Team lanci l\textquotesingle attacco, chiedete allo
studente che fa la vittima di aprire il prompt dei comandi e
visualizzare la propria cache ARP "pulita":

\begin{itemize}
\item
  \textbf{Comando (Windows/Linux):} arp -a
\item
  \emph{Fate annotare il MAC address associato all\textquotesingle IP
  del Default Gateway.}
\end{itemize}

Una volta che l\textquotesingle attaccante ha avviato arpspoof, fate
ripetere il comando arp -a alla vittima. Gli studenti dovranno notare
l\textquotesingle evidenza dell\textquotesingle avvelenamento:
\textbf{L\textquotesingle indirizzo IP del Gateway e
l\textquotesingle indirizzo IP della macchina Kali (attaccante)
presenteranno ora l\textquotesingle identico MAC Address.} Questo è il
segnale inequivocabile che il traffico sta venendo dirottato.

\paragraph{Step 2: La Mitigazione Locale (Tabelle ARP
Statiche)}\label{step-2-la-mitigazione-locale-tabelle-arp-statiche}

Per difendersi dall\textquotesingle attacco senza toccare le
configurazioni dello switch, la vittima può forzare il proprio sistema
operativo a ignorare i pacchetti ARP dinamici ingannevoli, inserendo una
regola statica (hardcoded) per il Default Gateway.

Fate pulire la cache inquinata e inserite la regola statica (sono
necessari i privilegi di Amministratore/Root):

\textbf{Su Windows (Terminale come Amministratore):}

DOS

arp -d *

arp -s \textless IP\_DEL\_GATEWAY\textgreater{}
\textless MAC\_REALE\_DEL\_GATEWAY\textgreater{}

\emph{(Nota per i docenti: Su Windows 10/11, il comando arp classico
potrebbe essere deprecato a favore di netsh. In tal caso usare: netsh
interface ipv4 add neighbors "Nome Scheda"
\textless IP\_GATEWAY\textgreater{} \textless MAC\_REALE\textgreater)}

\textbf{Su Linux/Kali:}

Bash

sudo ip -s -s neigh flush all

sudo arp -s \textless IP\_DEL\_GATEWAY\textgreater{}
\textless MAC\_REALE\_DEL\_GATEWAY\textgreater{}

\paragraph{Step di Verifica Finale}\label{step-di-verifica-finale}

Con la voce statica inserita, se la vittima lancia un arp -a, vedrà che
il tipo di record per il gateway è passato da "dinamico" a "statico".

Se a questo punto l\textquotesingle attaccante riprova a lanciare
arpspoof, il sistema operativo della vittima scarterà semplicemente i
pacchetti falsificati, mantenendo il traffico saldamente indirizzato
verso il MAC address reale del router. L\textquotesingle attacco è
neutralizzato.

\subsection{Denial of Service a Layer 2 - MAC
Flooding}\label{denial-of-service-a-layer-2---mac-flooding}

Mentre gli attacchi precedenti miravano a intercettare o dirottare il
traffico, il \textbf{MAC Flooding} punta a mandare in crisi
l\textquotesingle hardware dello switch, forzandolo a degradare le sue
prestazioni di sicurezza.

\subsubsection{Teoria: La CAM Table (Content Addressable
Memory)}\label{teoria-la-cam-table-content-addressable-memory}

Ogni switch mantiene una tabella interna chiamata \textbf{CAM Table} (o
\emph{MAC Address Table}). Questa tabella associa ogni indirizzo MAC
alla porta fisica a cui è collegato.

\begin{itemize}
\item
  \textbf{Il limite:} La memoria dello switch è finita (es. un Cisco
  2960 può contenere circa 8.000 indirizzi).
\item
  \textbf{L\textquotesingle attacco:} Se un attaccante inonda lo switch
  con migliaia di frame Ethernet, ognuno con un MAC address sorgente
  casuale e falso, la CAM Table si riempie in pochi secondi.
\item
  \textbf{Fail-open mode:} Quando la tabella è piena, lo switch non sa
  più dove mandare i pacchetti. Per non interrompere il servizio, entra
  in modalità "fail-open" e inizia a comportarsi come un \textbf{Hub}:
  inoltra ogni pacchetto in arrivo su \textbf{tutte le porte}
  (flooding).
\end{itemize}

\subsubsection{Laboratorio: Esecuzione con
macof}\label{laboratorio-esecuzione-con-macof}

Useremo \textbf{macof}, un tool della suite dsniff progettato
specificamente per generare rumore a Layer 2.

\begin{enumerate}
\def\labelenumi{\arabic{enumi}.}
\item
  \textbf{Lancio dell\textquotesingle attacco da Kali:} Apri il
  terminale e specifica l\textquotesingle interfaccia di rete:

  Bash

  \# Inizia a inondare lo switch con MAC address casuali

  sudo macof -i eth0

  \emph{Vedrai scorrere migliaia di righe con indirizzi IP e MAC
  generati casualmente.}
\item
  \textbf{Verifica sullo Switch (Visualizzare la CAM):} Mentre
  l\textquotesingle attacco è in corso, entra nella console dello switch
  Cisco per osservare il disastro in tempo reale.

  Cisco CLI

  \# Visualizza l\textquotesingle intera tabella dei MAC address

  Switch\# show mac address-table

  \# Per vedere quanti indirizzi sono stati memorizzati (utile per
  vedere il limite raggiunto)

  Switch\# show mac address-table count

  \emph{Noterai che la tabella si riempirà istantaneamente con MAC
  address senza senso, tutti associati alla tua porta (es. Gi0/1).}
\end{enumerate}

\subsubsection{Analisi dei Risultati e
Mitigazione}\label{analisi-dei-risultati-e-mitigazione}

Una volta che la CAM è piena, l\textquotesingle attaccante può
semplicemente lanciare \textbf{Wireshark} e vedrà transitare tutto il
traffico della rete (anche quello non destinato a lui), proprio come se
fosse collegato a un vecchio hub.

\paragraph{Come bloccare il MAC
Flooding?}\label{come-bloccare-il-mac-flooding}

La soluzione definitiva è la \textbf{Port Security}. Questa funzione
permette di limitare il numero di MAC address che possono essere appresi
da una singola porta.

\textbf{Comandi di Hardening sul 2960:}

Cisco CLI

Switch(config)\# interface range fa0/1 - 24

Switch(config-if)\# switchport mode access

Switch(config-if)\# switchport port-security

Switch(config-if)\# switchport port-security maximum 2

Switch(config-if)\# switchport port-security violation shutdown

\emph{Con questa configurazione, se macof prova a inviare più di 2 MAC
address differenti, lo switch spegne immediatamente la porta
(err-disable), bloccando l\textquotesingle attaccante.}

\subsubsection{Dynamic Host Configuration Protocol -
DHCP}\label{dynamic-host-configuration-protocol---dhcp}

Il Dynamic Host Configuration Protocol (DHCP) è uno dei servizi più
critici e, paradossalmente, meno protetti di una rete locale. Essendo un
protocollo basato su broadcast e privo di autenticazione, si presta a
diverse tipologie di attacco che possono portare alla paralisi della
rete (DoS) o all\textquotesingle intercettazione totale del traffico
(MitM).

\subsubsection{4.1 Tipologie di Attacchi
DHCP}\label{tipologie-di-attacchi-dhcp}

In ambito VA/PT, le principali minacce che analizzeremo sono:

\begin{enumerate}
\def\labelenumi{\arabic{enumi}.}
\item
  \textbf{DHCP Starvation:} Un attacco di Denial of Service che
  esaurisce tutti gli indirizzi IP disponibili nel pool del server
  legittimo.
\item
  \textbf{Rogue DHCP Server:} L\textquotesingle introduzione di un
  server non autorizzato che distribuisce parametri di rete malevoli
  (Default Gateway e DNS controllati dall\textquotesingle attaccante).
\item
  \textbf{DHCP Spoofing/Release:} L\textquotesingle invio di messaggi
  falsificati per terminare forzatamente i lease degli altri client o
  per intercettare i messaggi di acknowledge.
\end{enumerate}

\subsubsection{4.2 DHCP Starvation: Teoria e
Meccanica}\label{dhcp-starvation-teoria-e-meccanica}

L\textquotesingle attacco di \textbf{Starvation} (fame/esaurimento)
sfrutta la natura "generosa" del server DHCP. Quando un client entra in
rete, invia un pacchetto \emph{DHCP Discover} in broadcast. Il server
risponde con una \emph{DHCP Offer} riservando un IP per quel MAC
address.

L\textquotesingle attaccante utilizza un tool per inondare la rete con
migliaia di pacchetti \emph{DHCP Discover}, ognuno contenente un
\textbf{MAC address sorgente differente e casuale}. Il server DHCP (che
sia un router 4321 o un server Windows) esaurirà in pochi secondi
l\textquotesingle intero spazio di indirizzamento disponibile nel pool.

\textbf{Risultato:} Qualsiasi nuovo dispositivo legittimo che tenterà di
collegarsi alla rete rimarrà in stato di "Connettività limitata o
assente" poiché non riceverà alcuna risposta dal server.

\subsubsection{4.3 Laboratorio: Esecuzione con
Yersinia}\label{laboratorio-esecuzione-con-yersinia}

In questo scenario, utilizzeremo un \textbf{Cisco ISR 4321} come Server
DHCP e una macchina \textbf{Kali Linux} come attaccante.

\paragraph{1. Setup del Server (Cisco ISR
4321)}\label{setup-del-server-cisco-isr-4321}

Configurate il router per gestire il pool di indirizzi:

Cisco CLI

Router(config)\# ip dhcp pool LAN\_CLIENTS

Router(dhcp-config)\# network 10.10.10.0 255.255.255.0

Router(dhcp-config)\# default-router 10.10.10.1

Router(dhcp-config)\# lease 0 1 ! Lease di un\textquotesingle ora per il
test

\paragraph{2. Esecuzione dell\textquotesingle attacco (Kali
Linux)}\label{esecuzione-dellattacco-kali-linux}

L\textquotesingle attaccante avvia Yersinia in modalità interattiva:

\begin{enumerate}
\def\labelenumi{\arabic{enumi}.}
\item
  \textbf{Avvio:} sudo yersinia -I
\item
  \textbf{Selezione Interfaccia:} Premi i, seleziona
  l\textquotesingle interfaccia corretta e premi Invio.
\item
  \textbf{Cambio Protocollo:} Premi g, seleziona \textbf{DHCP} e premi
  Invio.
\item
  \textbf{Lancio Attacco:}

  \begin{itemize}
  \item
    Premi x per aprire il pannello degli attacchi.
  \item
    Seleziona 1 - sending RAW DHCP discover packets.
  \item
    Premi Invio.
  \end{itemize}
\item
  \textbf{Verifica:} Osserva il contatore dei pacchetti inviati nella
  schermata principale. Dopo circa 20 secondi, ferma
  l\textquotesingle attacco premendo L, selezionando
  l\textquotesingle attacco e cancellandolo (o chiudendo Yersinia con
  q).
\end{enumerate}

\paragraph{3. Verifica sul Router}\label{verifica-sul-router}

Eseguite il comando per vedere come il pool è stato saturato:

Cisco CLI

Router\# show ip dhcp binding

! Vedrete centinaia di righe con MAC address casuali

Router\# show ip dhcp server statistics

! Il contatore delle "Active Bindings" dovrebbe corrispondere alla
dimensione del pool

\subsubsection{Blue Team: Strategie di
Mitigazione}\label{blue-team-strategie-di-mitigazione}

La difesa contro la Starvation deve essere implementata sullo
\textbf{Switch di accesso (Cisco 2960)}, poiché è
l\textquotesingle apparato che vede fisicamente il traffico dei client.

\paragraph{1. DHCP Snooping}\label{dhcp-snooping}

Questa è la difesa principale. Lo switch monitora i messaggi DHCP e
permette i pacchetti di "risposta" (Offer/Ack) solo dalle porte fidate.

Cisco CLI

Switch(config)\# ip dhcp snooping

Switch(config)\# ip dhcp snooping vlan 10

! Definiamo la porta verso il router 4321 come TRUSTED

Switch(config)\# interface Gi0/1

Switch(config-if)\# ip dhcp snooping trust

\paragraph{2. Rate Limiting (Protezione specifica dalla
Starvation)}\label{rate-limiting-protezione-specifica-dalla-starvation}

Per bloccare l\textquotesingle inondazione di pacchetti \emph{Discover}
tipica di Yersinia, limitiamo il numero di pacchetti DHCP al secondo su
ogni porta client.

Cisco CLI

Switch(config)\# interface range fa0/1 - 24

Switch(config-if)\# ip dhcp snooping limit rate 10

\emph{Se un attaccante supera i 10 pacchetti DHCP al secondo, lo switch
metterà la porta in stato di \textbf{err-disable}, bloccando
l\textquotesingle attacco sul nascere.}

\paragraph{3. Port Security}\label{port-security}

Poiché la Starvation di Yersinia usa MAC address falsificati
all\textquotesingle interno del pacchetto DHCP, il \textbf{Port
Security} (visto nel Cap. 3) è un alleato fondamentale: se lo switch
vede troppi MAC sorgenti diversi provenire dalla stessa porta fisica,
bloccherà l\textquotesingle interfaccia prima ancora che il server DHCP
venga saturato.

In un attacco di Rogue DHCP, l\textquotesingle attaccante non è
obbligato a rispettare il piano di indirizzamento IP legittimo della
rete. Anzi, in un contesto di Penetration Testing professionale, spesso
è preferibile "spostare" la vittima su una \textbf{Shadow Network}: una
sottorete completamente diversa da quella ufficiale (es. la vittima si
aspetta la 10.10.10.0/24 e noi le assegniamo la 172.16.1.0/24).

\paragraph{Perché utilizzare una Subnet
diversa?}\label{perchuxe9-utilizzare-una-subnet-diversa}

\begin{enumerate}
\def\labelenumi{\arabic{enumi}.}
\item
  \textbf{Eliminazione dei conflitti IP:} Se usassimo la stessa subnet
  del router legittimo, rischieremmo di assegnare un IP già in uso,
  causando alert di sistema sulla vittima o collisioni ARP che
  renderebbero l\textquotesingle attacco instabile e rumoroso.
\item
  \textbf{Isolamento Logico:} Spostando la vittima su una rete che
  esiste solo nella memoria della macchina Kali,
  l\textquotesingle attaccante diventa l\textquotesingle unica autorità
  di routing. Il traffico della vittima non verrà mai intercettato o
  "corretto" dai gateway legittimi (come l\textquotesingle ISR 4321 o
  l\textquotesingle8000v).
\item
  \textbf{Trasparenza del MitM:} Attraverso il NAT (Network Address
  Translation), l\textquotesingle attaccante può far navigare la vittima
  verso internet facendole credere che tutto sia normale, mentre ogni
  singolo pacchetto viene decifrato e analizzato localmente.
\end{enumerate}

\subsubsection{Laboratorio: Configurazione della Shadow Network su
Kali}\label{laboratorio-configurazione-della-shadow-network-su-kali}

Per rendere efficace questo attacco, la macchina Kali deve essere pronta
a gestire una rete che ufficialmente non esiste sullo switch.

\paragraph{Step 1: Alias IP sull\textquotesingle interfaccia di
Kali}\label{step-1-alias-ip-sullinterfaccia-di-kali}

Affinché Kali possa comunicare con la vittima sulla nuova rete (es.
172.16.1.0/24), deve possedere un IP in quella subnet. Invece di
cambiare l\textquotesingle IP principale, creiamo un alias:

Bash

\# Sostituisci eth0 con la tua interfaccia bridge

sudo ip addr add 172.16.1.1/24 dev eth0

\paragraph{Step 2: Abilitazione del NAT
(Masquerading)}\label{step-2-abilitazione-del-nat-masquerading}

Per permettere alla vittima di uscire su internet tramite Kali (che
userà la sua connessione legittima), dobbiamo configurare il kernel
Linux affinché faccia da router NAT:

Bash

\# 1. Abilita l\textquotesingle IP Forwarding (se non già fatto)

sudo sysctl -w net.ipv4.ip\_forward=1

\# 2. Configura IPTables per il masquerading

\# Assumi che eth0 sia l\textquotesingle interfaccia verso la rete del
lab

sudo iptables -t nat -A POSTROUTING -o eth0 -j MASQUERADE

sudo iptables -A FORWARD -i eth0 -j ACCEPT

\paragraph{Step 3: Lancio del Rogue Server con
Yersinia}\label{step-3-lancio-del-rogue-server-con-yersinia}

Configurate Yersinia (tasto x -\textgreater{} opzione 2) con i parametri
della Shadow Network:

\begin{itemize}
\item
  \textbf{Pool Start/End:} 172.16.1.100 - 172.16.1.200
\item
  \textbf{Subnet Mask:} 255.255.255.0
\item
  \textbf{Default Gateway:} 172.16.1.1 (L\textquotesingle IP di Kali)
\item
  \textbf{DNS Server:} 172.16.1.1 (O un DNS pubblico come 8.8.8.8)
\end{itemize}

\subsubsection{Verifica dell\textquotesingle Attacco
(Post-Exploitation)}\label{verifica-dellattacco-post-exploitation}

Una volta che la vittima ha ottenuto l\textquotesingle IP dalla Shadow
Network, lo studente deve verificare la totale sottomissione del
traffico:

\begin{enumerate}
\def\labelenumi{\arabic{enumi}.}
\item
  \textbf{Analisi sulla Vittima:}

  \begin{itemize}
  \item
    ipconfig /all (Windows) o ip a (Linux): Verificare che
    l\textquotesingle IP appartenga alla rete 172.16.1.x.
  \item
    tracert 8.8.8.8: Il primo salto deve essere obbligatoriamente
    l\textquotesingle IP di Kali (172.16.1.1).
  \end{itemize}
\item
  \textbf{Analisi sull\textquotesingle Attaccante (Kali):}

  \begin{itemize}
  \item
    Eseguire tcpdump -i eth0 icmp e chiedere alla vittima di fare un
    ping.
  \item
    \textbf{Risultato atteso:} Vedrete passare i pacchetti della vittima
    attraverso Kali. Ora potete lanciare tool di sniffing come
    \textbf{Wireshark} o \textbf{Bettercap} per catturare password o
    sessioni HTTP.
  \end{itemize}
\end{enumerate}

\begin{quote}
\textbf{Nota di Troubleshooting:} Se la vittima riceve ancora
l\textquotesingle IP dal router legittimo, significa che il router è
stato più veloce di Kali. In questo caso, è necessario lanciare un
attacco di \textbf{DHCP Starvation} (visto nel paragrafo 4.2) per
"zittire" il router prima di riprovare con il Rogue Server.
\end{quote}

\subsubsection{DNS Spoofing: Il controllo nominale tramite
DHCP}\label{dns-spoofing-il-controllo-nominale-tramite-dhcp}

Mentre nel Capitolo 3 abbiamo visto come il DNS Spoofing richieda una
tecnica di intercettazione attiva (ARP Poisoning),
l\textquotesingle attacco tramite \textbf{Rogue DHCP Server} rende
questa fase estremamente più semplice e "pulita".

Poiché l\textquotesingle attaccante ha il controllo del campo
\textbf{Option 6 (Domain Name Server)} nel pacchetto di risposta DHCP,
può istruire la vittima a utilizzare la macchina Kali come unico server
DNS della rete. In questo scenario, non c\textquotesingle è bisogno di
intercettare query dirette a server terzi: la vittima invierà
\emph{volontariamente} ogni sua richiesta di risoluzione
all\textquotesingle attaccante.

\textbf{La catena di compromissione:}

\begin{enumerate}
\def\labelenumi{\arabic{enumi}.}
\item
  \textbf{Rilevamento:} La vittima si connette e riceve un IP dal Rogue
  Server.
\item
  \textbf{Configurazione:} Il resolver della vittima viene impostato
  sull\textquotesingle IP di Kali (es. 172.16.1.1).
\item
  \textbf{Risoluzione:} Quando l\textquotesingle utente cerca
  www.bancapopolare.it, la query UDP 53 arriva direttamente a Kali.
\item
  \textbf{Spoofing:} Un tool sull\textquotesingle attaccante (come
  dnschef o bettercap) intercetta la richiesta e risponde con
  l\textquotesingle IP di un server malevolo o di una pagina di
  phishing, inoltrando invece le richieste "non interessanti" ai veri
  DNS di Google (8.8.8.8) per non destare sospetti.
\end{enumerate}

\textbf{Vantaggio tattico:} A differenza dell\textquotesingle ARP
Spoofing, che può causare instabilità di rete o essere rilevato da
sistemi IDS (come l\textquotesingle Arpwatch), il DNS Spoofing via DHCP
è estremamente silenzioso a livello di protocollo, poiché si basa su
un\textquotesingle interazione client-server perfettamente standard. La
vittima non vedrà mai "due MAC address per lo stesso IP", ma
semplicemente un server DNS che sta facendo il suo lavoro.

\subsubsection{Sicurezza delle VLAN e VLAN
Hopping}\label{sicurezza-delle-vlan-e-vlan-hopping}

In una rete aziendale, le VLAN servono a isolare il traffico (es. VLAN
Guest separata dalla VLAN Server). Tuttavia, la segmentazione non è
magica: se l\textquotesingle infrastruttura non è configurata
correttamente, un attaccante può "saltare" i confini delle VLAN senza
passare per un firewall. Questo è il \textbf{VLAN Hopping}.

\subsubsection{Anatomia dell\textquotesingle Attacco: Il "Trunk
Misconfiguration"}\label{anatomia-dellattacco-il-trunk-misconfiguration}

Esistono due modi principali per eseguire il VLAN Hopping:

\begin{enumerate}
\def\labelenumi{\arabic{enumi}.}
\item
  \textbf{Switch Spoofing (Abuso del DTP):} L\textquotesingle attaccante
  usa Yersinia per negoziare automaticamente un trunk.
\item
  \textbf{Fixed Trunk (Errore sistemistico):}
  L\textquotesingle amministratore lascia accidentalmente una porta host
  configurata come trunk. In questo caso, lo switch non chiede chi sei:
  accetta qualsiasi pacchetto taggato gli arrivi.
\end{enumerate}

\subsubsection{Laboratorio: Infiltrazione e
Post-Exploitation}\label{laboratorio-infiltrazione-e-post-exploitation}

In questo scenario, simuliamo una porta switch (Gi0/1) erroneamente
configurata come Trunk. Entreremo nella \textbf{VLAN 10 (Segreta)} e
useremo Yersinia per attaccare i servizi interni.

\paragraph{1. Setup dello Switch
(vIOS-L2)}\label{setup-dello-switch-vios-l2}

Creiamo la zona protetta e apriamo il "tunnel" per
l\textquotesingle attaccante:

Cisco CLI

Switch\# conf t

Switch(config)\# vlan 10

Switch(config-vlan)\# name AREA\_PROTETTA

Switch(config-vlan)\# exit

Switch(config)\# interface GigabitEthernet 0/1

Switch(config-if)\# switchport trunk encapsulation dot1q

Switch(config-if)\# switchport mode trunk

Switch(config-if)\# no shutdown

\paragraph{2. Infiltrazione da Kali Linux (VLAN
Hopping)}\label{infiltrazione-da-kali-linux-vlan-hopping}

Dobbiamo dire al kernel Linux di gestire i tag 802.1Q per "parlare" con
la VLAN 10.

Bash

\# Carichiamo il modulo per il tagging

sudo modprobe 8021q

\# Creiamo l\textquotesingle interfaccia virtuale legata alla VLAN 10

sudo ip link add link eth0 name eth0.10 type vlan id 10

sudo ip link set dev eth0.10 up

\# Assegniamo un IP per muoverci nella nuova rete

sudo ip addr add 10.10.10.100/24 dev eth0.10

\paragraph{3. Attacco con Yersinia
(Post-Exploitation)}\label{attacco-con-yersinia-post-exploitation}

Ora che Kali è "dentro" la VLAN 10, possiamo usare Yersinia per creare
il caos nel segmento protetto:

\begin{itemize}
\item
  \textbf{STP Attack:} Diventare il Root Bridge della VLAN 10 per
  intercettare tutto il traffico (Sniffing).

  \begin{itemize}
  \tightlist
  \item
    In Yersinia: g -\textgreater{} STP -\textgreater{} x -\textgreater{}
    1 (Claiming Root Role).
  \end{itemize}
\item
  \textbf{DHCP Starvation:} Esaurire gli IP della VLAN 10 per bloccare i
  nuovi client legittimi.

  \begin{itemize}
  \tightlist
  \item
    In Yersinia: g -\textgreater{} DHCP -\textgreater{} x
    -\textgreater{} 1 (Raw Discover packets).
  \end{itemize}
\end{itemize}

\subsubsection{Blue Team: Mitigazione e
Hardening}\label{blue-team-mitigazione-e-hardening}

Proteggere le VLAN è semplice, ma richiede disciplina. Ecco le tre
regole d\textquotesingle oro da applicare sul Catalyst 2960:

\begin{longtable}[]{@{}
  >{\raggedright\arraybackslash}p{(\linewidth - 4\tabcolsep) * \real{0.1948}}
  >{\raggedright\arraybackslash}p{(\linewidth - 4\tabcolsep) * \real{0.3247}}
  >{\raggedright\arraybackslash}p{(\linewidth - 4\tabcolsep) * \real{0.4545}}@{}}
\toprule\noalign{}
\begin{minipage}[b]{\linewidth}\raggedright
Azione
\end{minipage} & \begin{minipage}[b]{\linewidth}\raggedright
Comando Cisco
\end{minipage} & \begin{minipage}[b]{\linewidth}\raggedright
Perché farlo?
\end{minipage} \\
\midrule\noalign{}
\endhead
\bottomrule\noalign{}
\endlastfoot
\textbf{Forza modalità Access} & switchport mode access & Impedisce alla
porta di diventare un Trunk, anche se riceve pacchetti DTP. \\
\textbf{Disabilita DTP} & switchport nonegotiate & Spegne completamente
il protocollo di negoziazione. \\
\textbf{Cambia Native VLAN} & switchport trunk native vlan 999 &
Impedisce l\textquotesingle attacco di \textbf{Double Tagging} (usando
una VLAN inutilizzata). \\
\end{longtable}

Esporta in Fogli

\paragraph{Configurazione di Hardening
consigliata:}\label{configurazione-di-hardening-consigliata}

Cisco CLI

Switch(config)\# interface range fa0/1 - 24

Switch(config-if)\# switchport mode access

Switch(config-if)\# switchport nonegotiate

Switch(config-if)\# spanning-tree bpduguard enable

\end{document}
